\section{Códigos reportados como no funcionales}

¿Un código no te sirvió? Ayúdanos a mantener la lista limpia. Aquí se registran los códigos que ya no funcionan, para que entre todos los revisemos y los borremos de su archivo original.

\begin{myquote}{Cómo reportar un código}
\begin{enumerate}[leftmargin=*]
  \item Abre \href{\ghedit{reportes.tex}}{este archivo en GitHub} y edítalo.
  \item Agrega un reporte con el comando:
        \texttt{\textbackslash reporte\{CODIGO\}\{Carrera -- Plantel\}\{Motivo\}\{ruta/archivo.tex\}}
  \item Cada reporte genera dos enlaces al archivo original: \textbf{Ver archivo} y
        \textbf{Editar y borrar}, para que puedas revisarlo y eliminarlo directo en GitHub.
\end{enumerate}
\end{myquote}

\subsection*{Reportes abiertos}
\begin{itemize}[leftmargin=*]
  % Agrega tus reportes aqui, uno por linea, con:
  %   \reporte{CODIGO}{Carrera -- Plantel}{Motivo}{ruta/archivo.tex}
  \item \textit{Por ahora no hay reportes abiertos.}
\end{itemize}

\noindent\textit{Si ya resolviste un reporte (borraste el código del archivo original), puedes quitar también su entrada de esta lista.}